% Overview by GPT-6 Astra (OpenAI, ChatGPT Pro mode).
\begingroup
\selectlanguage{english}
\noindent\textbf{Reconstruction, preservation, and the geometry of indexing categories.}

These pages complete the second chapter by asking how much of a category is
visible to its diagram theories. The starting point is the cofinality theorem:
a restriction of diagrams preserves a specified limit or colimit when the
appropriate comma categories are connected. The next question is more demanding.
When does restriction embed the whole category of diagrams fully faithfully?
Grothendieck then turns from individual functors to classes of functors. He
tries to recover the classes controlling change of base from a chosen class
of weak equivalences, and develops the order-theoretic machinery needed to
state such a recovery. The closing discussion of cofiltered categories gives
concrete tests and counterexamples for this programme. These are successive
parts of one investigation: which local tests retain exactly the information
that a theory of coefficients can detect?

For ordinary presheaves, the answer to the full-faithfulness question is a
two-sided cofinality criterion. Given $u:A\to B$ and an arrow
$\gamma:b_0\to b_1$, form the category of factorizations
$b_0\to u(a)\to b_1$ whose composite is $\gamma$. Restriction on
set-valued presheaves is fully faithful precisely when every such category
is nonempty and connected. The counit evaluated on a representable presheaf
has these connected-component sets as its fibres, and representables generate
all presheaves by colimits. This is the decisive argument of Proposition 4.1.
In later terminology the functors are lax epimorphisms, or absolutely dense
functors; the equivalence of these formulations is established by
\citet{AdamekElBashirSobralVelebil2001}. The overlap with the preceding package
already treats the general criterion. Here the commentary supplies the
missing distinction between its two adjunction tests and checks the direction
of the comma-category maps.

The right-adjoint unit behaves differently. Its being invertible on every
representable says that the restricted Yoneda functor
$B\to\operatorname{Set}^{A^{\mathrm{op}}}$ is fully faithful: this is
density, in the usual sense \citep[Definition 6.5.11]{Riehl2016}. It does not
imply full faithfulness of restriction on all presheaves. The full category
of sets on a singleton and a two-element set, with $u$ selecting the
singleton, already gives a counterexample. All maps of the original category
are recovered from maps of sets of elements, but the identity of the
two-element set cannot factor through the singleton. Thus the hesitation on
manuscript pages 54--56 concerns a genuine obstruction. The commentary also
computes the four-object example on pages 57--58: every factorization
category is contractible, while one strict pullback is a disconnected
two-point category. The example separates reconstruction of diagrams from
universal preservation under strict base change.

For an arbitrary derivator $\D$, one must distinguish full faithfulness for
that particular coefficient theory from full faithfulness for every
derivator. Modern homotopical epimorphisms impose the latter quantifier.
\citet[Definition 6.1 and Theorem 6.10, arXiv v3]{GrothStovicek2016Trees}
characterize them by weak contractibility of the same factorization
categories. Grothendieck's fixed-$\D$ problem is subtler. Global diagrams
test only certain diagrams on a slice, and his proof of Proposition 4.4
carefully isolates the additional detection needed to recover all the slice
tests. A conservative left Kan extension from the slice suffices; essential
surjectivity of restriction to the slice is a stronger sufficient condition.
The pointwise Kan-extension and detection axioms of a derivator perform
different jobs in this argument \citep[Definition 1.5]{Groth2013}.

The developments in the commentary make two limits of these sufficient
conditions explicit. First, in the derivator represented by the complete
lattice $\{0<1\}$, diagrams on the free category generated by a directed
two-cycle are constant. Evaluation at one vertex is therefore fully
faithful. Nevertheless, the identity of the other vertex has no
factorization through it, and its factorization category is empty. This
answers the intended question about $\D$-connectedness versus
$\D$-biasphericity. Second, the parenthetical faithfulness assertion in
Corollary 4.7 fails for coherent derived diagrams. A discrete fibration from
a three-object fork to an interval sends a nonsplit extension of quiver
representations to a split extension. Direct sums in the coefficient
derived category are faithful, but the Kan extension is not. The
conservativity assertion survives: pointwise detection proves it. Each
development is given with a complete proof and a statement of the sources
checked, without a claim of priority.

The cancelled passage at the end of \S4 must be read with comparable care.
Let $S(A)$ be the maps whose pullback functors preserve a class $A$, and let
$N(L)$ be the maps whose pullbacks along a testing class $L$ lie in $W$.
Under the requisite identity and stability hypotheses, the inclusions and
the two idempotence calculations in the passage survive. The missing
ingredient is order reversal of $S$: enlarging the preserved class changes
both the inputs and the required outputs of a pullback test. It need not
give a smaller stabilizer. With $W$ the maximal localizer on categories and
$A$ the categorical equivalences, two point inclusions into the walking
isomorphism directly refute the stronger identity used in the printed
justification. This diagnosis preserves the valid calculations while
explaining why they do not constitute a Galois connection.

Section 5 supplies that connection in its proper generality. Any relation
between two sets defines two antitone polar operations. The composites are
closure operators, and their closed subsets form anti-isomorphic complete
lattices. When the sets are ordered, a further condition makes polar lower
sets principal and hence represents their closures by elements. Formula
(5.20) slips between these levels: a closed subset is not its greatest
element, and a join of closed elements need not be closed in the ambient
lattice. The corrected formula takes the closure of the join. Double
annihilation in $\mathbf F_2^2$ gives a short counterexample to the printed
equality. The classical examples also require their proper separation or
topological hypotheses. In infinite Galois theory the fixed subgroups are
the Krull-closed subgroups; the dense cyclic Frobenius subgroup in
$\widehat{\mathbf Z}$ shows why the qualification cannot be dropped
\citep[Proposition 7.12 and Theorem 7.13]{Milne2022FT}. A lifting relation is
another polarity, but a mutually closed pair of lifting classes still
requires a separate factorization theorem to become a weak factorization
system \citep[Definition 11.2.1]{Riehl2014}.

Section 6 repairs the earlier construction by fixing its target class.
The transporter $T_W(A)$ consists of maps whose pullback functors send
arrows of $A$ to arrows of $W$. Testing includes every arrow in the slice,
with any structure map from its codomain to the base. It is not enough to
test arrows whose codomain happens to equal the base. With $A$ closed under
composition and $L$ stable under pullback, both containing isomorphisms,
the equivalence
\[
 A\subseteq N_W(L)\quad\Longleftrightarrow\quad L\subseteq T_W(A)
\]
is a genuine antitone adjunction. The correction at (6.7) removes
composition closure from the right-hand domain, where it is not automatic.
A four-object chain in the commentary exhibits a closed transporter that
does not compose, under exactly the abstract assumptions of \S6. The
example is not claimed to use a fundamental localizer on $\mathrm{Cat}$.
Proposition 6.1 then identifies the additional condition: for a closed pair,
the right class is composition-closed exactly when it equals the stabilizer
of its left class.

This formalism has a direct historical continuation. Grothendieck's letter
to Thomason singles out a fibration class defined intrinsically from weak
equivalences, while treating factorization as another question
\citep[\S1, pp.~3--4]{Grothendieck1991Letter}. Rezk's sharp maps use the same
preservation condition. In a right proper model category they characterize
the maps all of whose strict pullback squares are homotopy cartesian
\citep[Propositions 2.2 and 2.7]{Rezk1998Sharp}. Maltsiniotis develops local
tests for the corresponding preservation of aspherical maps
\citep[Definition 4.6 and Theorem 4.13]{Maltsiniotis2005Asphericite}.
One question on page 89 has a particularly concrete answer: the stabilizer
of categorical equivalences is exactly the isofibrations. The reader gives
both directions of the proof, using the walking isomorphism for necessity,
and relates the result to the categorical model structure
\citep[\S2 and Theorem 3.1]{Rezk1996Categories}. For the general mixed
transporters and the proposed bilocal description of the biaspherical
stabilizer, the report states the extent of verification. Failure to find
a precise matching theorem is not presented as proof that the questions
are open today.

The final section turns the abstract machinery into the geometry of
predecessors. A cofiltered functor is tested by cofiltered comma categories;
the local version allows several components and asks each to be
cofiltered. Over a poset, these conditions become statements about lower
open subsets. Proposition 7.4 is essentially sound, but its target lower
open must be generated by the image of the chosen open subset: taking the
whole principal lower set makes the printed condition false even for the
empty subset. The commentary proves the corrected criterion and the
initial-object assertion of Remark 7.5. It also exhibits a cofiltered map
whose arbitrary base change is not even locally weakly cofiltered, while
respecting the stated preservation under co-étale base change. For
fundamental localizers, the later minimality theorem removes the extra
Loc(8) qualification attached here to the asphericity of cofiltered
categories \citep[Theorem 2.2.11]{Cisinski2004Minimal}.

The final diagrams distinguish common predecessors, the filling of
cospans, equalization of parallel arrows, and connectedness of the
categories of cones over every finite family of objects. Their implications
become transparent when these requirements are
kept separate. Groups, the simplex category, a parallel-arrow category,
and disjoint unions give the required separating examples. The modern
interpretation of total connectedness is siftedness of the opposite
category; cofilteredness adds the equalization needed for commuting with
all finite limits, rather than only finite products
\citep{AdamekRosickyVitale2010,StacksFiltered2026,Tendas2024}. In this way the
chapter ends where its larger project began: the shape of an indexing
category governs which constructions a theory of diagrams can perform
and which information restriction can retain. The derived-diagram
motivation already appears in \emph{Pursuing Stacks}, \S69
\citep{Grothendieck1983}; here it is being refined into local, verifiable
conditions on functors.

\medskip
\noindent\emph{Commentary and audit by GPT-6 Astra (OpenAI, ChatGPT Pro mode).
This credit concerns the work supplied here.}
\nocite{*}
\endgroup
